\section{Agente SAM3: arquitectura de colaboración entre VLM y SAM3}

\subsection{Diseño de la arquitectura: cerebro y ojos/manos}

El repositorio oficial de SAM3 ofrece un ejemplo de agente cuya idea central es dividir el trabajo entre el VLM y SAM3:

\begin{importantbox}{Modo de división del trabajo del agente SAM3}
\begin{itemize}
    \item \textbf{VLM (cerebro)}: encargado del razonamiento semántico de alto nivel, de entender la intención en lenguaje natural del usuario y de planear los pasos de ejecución
    \item \textbf{SAM3 (ojos/manos)}: encargado de la operación visual de bajo nivel, de generar segmentaciones de Mask precisas
\end{itemize}
El VLM no dibuja el Mask directamente, sino que le indica a SAM3 que lo dibuje.
\end{importantbox}

\subsection{Definición de herramientas y System Prompt}

El agente SAM3 define cuatro herramientas (Tools), de las cuales la más central es \texttt{segment\_phrase}:

\begin{itemize}
    \item Entrada: una frase nominal simple (como "child", "blue vest")
    \item Salida: la imagen de Mask de todos los objetos coincidentes en la imagen completa (con etiquetas numeradas)
\end{itemize}

El System Prompt establece explícitamente restricciones clave:
\begin{itemize}
    \item \textbf{Prohibido usar palabras de posición}: no se puede escribir "left child"
    \item \textbf{Prohibido usar cuantificadores}: no se puede escribir "three cats"
    \item \textbf{Sustantivos genéricos obligatorios}: solo se pueden usar frases genéricas como "child", "person"
\end{itemize}

\subsection{Flujo de interacción completo de Visual Grounding}

Tomando como ejemplo "encontrar al niño más a la izquierda que trae puesto un chaleco azul":

\textbf{Ronda 1: etapa de segmentación}:
\begin{enumerate}
    \item GPT-4o recibe la imagen original y la Query del usuario
    \item Como SAM3 no es hábil para el razonamiento de posición (esto ya se indica en la descripción de la Tool), GPT-4o planea una estrategia: primero segmentar a todos los niños que traen puesto un chaleco azul
    \item Se invoca \texttt{segment\_phrase("child in blue vest")}
    \item SAM3 devuelve un diagrama estructurado con Mask numerados (se detectaron 3 instancias coincidentes)
\end{enumerate}

\textbf{Ronda 2: etapa de razonamiento}:
\begin{enumerate}
    \item GPT-4o recibe el diagrama estructurado anotado con los Mask (¡esto es precisamente Visual Prompting!)
    \item Realiza verificación visual y razonamiento espacial para determinar qué Mask ID corresponde al niño "más a la izquierda"
    \item Emite la respuesta final
\end{enumerate}

\begin{warningbox}{El razonamiento espacial del VLM aún tiene defectos}
En una prueba real, GPT-4o determinó que el Mask ID 1 era el más a la izquierda, pero la respuesta correcta debía ser el Mask ID 2: el modelo confundió izquierda y derecha. Esto muestra que incluso el VLM más avanzado sigue teniendo inestabilidad en el razonamiento de posición espacial.
\end{warningbox}

\subsection{Detalles de la implementación}

La implementación oficial del agente SAM3 no usa marcos de agentes modernos como LangChain o LangGraph, sino que adopta un enfoque más simple:

\begin{itemize}
    \item Admite entrada mixta de texto e imagen
    \item La definición de la Tool se escribe directamente en el System Prompt
    \item El ejemplo oficial se basa en LLaMA 3 Vision 8B; el ponente lo cambió por GPT-4o (que admite Tool y entrada multimodal)
\end{itemize}

\subsection{Resumen de la sección}

El agente SAM3 muestra un paradigma clásico de colaboración entre "VLM y un modelo visual especializado": el VLM se encarga del razonamiento y la planeación, y SAM3 se encarga de la operación visual precisa. Mediante el mecanismo de Tool Call, un VLM que no cuenta con capacidad de segmentación adquiere una capacidad de segmentación a nivel de píxel, con lo que se logra el Visual Grounding.
